% TOP_PROBLEM: {"id": 2306070, "title": "Research Problems in Function Theory — Problem 6.70", "category_id": 8, "category": {"id": 8, "name": "analysis", "display_name": "Analysis", "description": "Limits, continuity, calculus, and function theory.", "slug": "analysis", "order_index": 8, "created_at": "2026-07-31T15:26:25.671Z"}, "status": "open"}
% FIRST_POSTED: null
% ENTRY_KIND: statement_only
% Imported from: batch03.tex; UnsolvedMath ID 2306070
\documentclass[11pt]{article}
\usepackage[margin=25mm]{geometry}
\usepackage{fontspec}
\setmainfont{Latin Modern Roman}
\usepackage{amsmath,amssymb,mathtools}
\usepackage{unicode-math}
\setmathfont{Latin Modern Math}
\usepackage{xurl,hyperref}
\hypersetup{colorlinks=true,urlcolor=blue}
\setcounter{secnumdepth}{0}
\setlength{\parindent}{0pt}
\setlength{\parskip}{5pt}
\setlength{\emergencystretch}{3em}
\newcommand{\sref}[2]{\hyperlink{source:#1:#2}{[S#2]}}
\newcommand{\cataloguescope}[1]{\paragraph{Recorded target.}#1}
\begin{document}
% UnsolvedMath ID 2306070; source catalogue code AMR-022-6070
\setcounter{section}{2306069}
\section{Research Problems in Function Theory — Problem 6.70}\label{2306070}
{\small\sffamily UnsolvedMath ID 2306070 · Analysis}
\subsection{Definitions and mathematical statement}

\paragraph{Shared notation.}$\mathbb N=\{1,2,\ldots\}$, and $\mathbb Z,\mathbb Q,\mathbb R,\mathbb C$ denote the integers, rationals, reals and complex numbers. Any other conventions are specified locally.


Write $\mathbb D=\{z\in\mathbb C:|z|<1\}$ and $\mathbb T=\partial\mathbb D$. The class $S$ consists of holomorphic injective functions $f:\mathbb D\to\mathbb C$ normalized by $f(0)=0$ and $f'(0)=1$, so $f(z)=z+\sum_{n=2}^{\infty}a_nz^n$. A function $f$ in a class $\mathcal A$ is an extreme point of that class if $f=tg+(1-t)h$, with $g,h\in\mathcal A$ and $0<t<1$, forces $g=h=f$. Equip $H(\mathbb D)$, the vector space of holomorphic functions on $\mathbb D$, with uniform convergence on compact subsets. A support point of $S$ maximizes $\operatorname{Re}L$ over $S$ for a continuous complex-linear functional $L:H(\mathbb D)\to\mathbb C$ whose real part is not constant on $S$.  This class-relative extreme-point convention and the distinct closed-convex-hull convention are stated separately in \sref{2306070}{3}.
\paragraph{Statement recorded in the supplied source.}
Is every extreme point of $S$ a support point? Conversely, is every support point of $S$ an extreme point? \sref{2306070}{2}
Update 6.70 records counterexamples to the first implication by Hamilton and by Duren and Leung; it reports the converse as unresolved at the time of the update. \sref{2306070}{2}
\subsection{Short English statement}
Do extreme points and support points coincide in the normalized univalent class?
\subsection{Sources}
\begin{description}
\item[\hypertarget{source:2306070:1}{S1}] ULAM.AI and the UnsolvedMath Contributors, \emph{UnsolvedMath: A Curated Collection of Open Mathematics Problems}, record 2306070 (AMR-022-6070). CC BY 4.0. \url{https://huggingface.co/datasets/ulamai/UnsolvedMath}.
\item[\hypertarget{source:2306070:2}{S2}] W. K. Hayman and E. F. Lingham, \emph{Research Problems in Function Theory} (2018), Chapter 6, Problem 6.70 and its complete update; Chapter 6 initial notation. \url{https://arxiv.org/abs/1809.07200}.
\item[\hypertarget{source:2306070:3}{S3}] V. Allu and A. Pandey, \emph{Support points of some classes of analytic and univalent functions} (2022), Introduction and Preliminaries, p.3. \url{https://arxiv.org/pdf/2209.10607}.
\end{description}
\label{end:2306070}
\end{document}
